\chapter{Report review, comparison and drafting}
\label{ch:review}

Three features share a theme: they look at reports as \emph{documents to be
checked}, not as a corpus to be searched.

\section{Report review}
\label{sec:report-review}

The review skill inspects a technical report before it reaches a human
reviewer. It does not re-solve the CAE analysis and it never certifies physical
correctness. What it does is check that the report says what a report of its
discipline is expected to say, and show the evidence for every complaint.

\subsection{Findings}

Every rule produces zero or more findings with a fixed shape.

\begin{longtable}{L{0.24\textwidth} L{0.66\textwidth}}
\toprule
\textbf{Field} & \textbf{Meaning} \\
\midrule
\endfirsthead
\toprule
\textbf{Field} & \textbf{Meaning} \\
\midrule
\endhead
\bottomrule
\endfoot
\code{rule\_id} & Stable identifier, e.g. \code{captions.sequence} or
\code{nvh.measurement\_setup}. \\
\code{category} & \code{structure}, \code{captions}, \code{numbers},
\code{extraction}, \code{content}, or the discipline name. \\
\code{severity} & \code{critical}, \code{warning} or \code{info}. \\
\code{status} & The engine's verdict, e.g. \code{needs\_review}. \\
\code{message} / \code{suggested\_fix} & What is wrong, and what to do about
it. \\
\code{evidence} & The quoted lines the finding rests on, page-prefixed. \\
\code{page\_start} / \code{page\_end} & Where in the report, used to drive the
highlighted PDF preview. \\
\code{engine} & \code{rules} or \code{semantic}. \\
\code{finding\_key} & A hash over the identifying parts of the finding. It is
what a human decision is recorded against, so the same finding on a re-run keeps
its decision. \\
\end{longtable}

\subsection{Deterministic rules}

These run with no model and apply to every report.

\begin{longtable}{L{0.30\textwidth} L{0.60\textwidth}}
\toprule
\textbf{Rule} & \textbf{Checks} \\
\midrule
\endfirsthead
\toprule
\textbf{Rule} & \textbf{Checks} \\
\midrule
\endhead
\bottomrule
\endfoot
\code{metadata.required\_fields} & Cover fields: report number, date, prepared
by, checked by. \\
\code{structure.required\_sections} & Presence of a scope section and a
results/conclusion section. \\
\code{captions.sequence} & Table and figure numbering runs in order with no gaps
or repeats. \\
\code{captions.title} & Every table and figure caption actually has a title. \\
\code{captions.references} & Captions are referenced from the body text, and
references resolve. \\
\code{numbers.decimal\_style} & Decimal separators are used consistently in
measured values. \\
\code{extraction.no\_text} & Pages whose text could not be read at all. \\
\code{extraction.ocr\_low\_quality} & Pages where OCR produced text of doubtful
quality. \\
\code{content.embedded\_paths} & Local or network file paths left in the report
body. \\
\end{longtable}

The last two categories are why extraction provenance is stored per page: a
report the engine cannot read is a finding about the \emph{report}, not a silent
gap in the review.

\subsection{Discipline profiles}

Discipline rules are \textbf{data, not code}. One JSON file per discipline lives
in \file{app/rules/profiles/}.

\begin{longtable}{L{0.20\textwidth} L{0.18\textwidth} L{0.50\textwidth}}
\toprule
\textbf{File} & \textbf{Label} & \textbf{Rules} \\
\midrule
\endfirsthead
\bottomrule
\endfoot
\file{nvh.json} & NVH & Measurement setup, signal processing, result
interpretation. \\
\file{cfd.json} & CFD & Solution model and boundary conditions, mesh,
convergence. \\
\file{durability.json} & Durability & Material definition, loads and supports,
mesh, contacts, result-to-criterion link. \\
\file{test.json} & Test / Validasyon & Test object and configuration, equipment,
conditions, procedure, acceptance decision, calibration. \\
\end{longtable}

Every discipline rule runs through one generic handler: each
\code{requirement\_groups} entry must be mentioned somewhere in the report by at
least one of the wordings it lists, and the groups that are not mentioned are
named in the finding message.

\subsection{The profile file schema}
\label{sec:profile-schema}

\begin{jsonblock}
{
  "profile": "nvh",
  "label": "NVH",
  "detect_priority": 10,
  "aliases": ["nvh", "noise vibration harshness"],
  "detect_patterns": ["\\bnvh\\b", "titresim"],
  "rules": [
    {
      "rule_id": "nvh.measurement_setup",
      "label": "NVH olcum duzeni ve kosullari",
      "category": "nvh",
      "severity": "warning",
      "message": "NVH olcum duzeni raporda tam izlenemiyor.",
      "suggested_fix": "Sensoru ve olcum noktasini, eksenleri ve her
                        olcumdeki calisma kosulunu yazin.",
      "requirement_groups": [
        {
          "label": "olcum noktasi / sensor",
          "aliases": ["sensor", "ivmeolcer", "akselerometre"]
        }
      ]
    }
  ]
}
\end{jsonblock}

Top-level keys --- all six required, no others accepted:

\begin{longtable}{L{0.22\textwidth} L{0.16\textwidth} L{0.52\textwidth}}
\toprule
\textbf{Key} & \textbf{Type} & \textbf{Meaning and constraints} \\
\midrule
\endfirsthead
\bottomrule
\endfoot
\code{profile} & string & Must equal the file name without \file{.json}.
\code{general} and \code{auto} are reserved for the built-in profiles. \\
\code{label} & string & The profile name shown in the interface. \\
\code{detect\_priority} & integer & Order in which \code{auto} tries the
identity patterns; lower goes first. A profile whose patterns are broad belongs
later in the sequence. \\
\code{aliases} & list of strings & Non-empty, no duplicates. What a catalog
discipline value may say to select this profile. \\
\code{detect\_patterns} & list of strings & Non-empty regular expressions, each
compiled at load time, matched against the report identity when no catalog
discipline settles it. \\
\code{rules} & list of objects & Non-empty. The rules below. \\
\end{longtable}

Each entry in \code{rules} --- again all seven keys required, no others
accepted:

\begin{longtable}{L{0.26\textwidth} L{0.15\textwidth} L{0.49\textwidth}}
\toprule
\textbf{Key} & \textbf{Type} & \textbf{Meaning and constraints} \\
\midrule
\endfirsthead
\toprule
\textbf{Key} & \textbf{Type} & \textbf{Meaning and constraints} \\
\midrule
\endhead
\bottomrule
\endfoot
\code{rule\_id} & string & Must start with \code{<profile>.}. That prefix is
what keeps ids unique across files without a global registry; a collision with a
built-in or semantic rule id is rejected. \\
\code{label} & string & Non-empty. The rule name in the interface. \\
\code{category} & string & Non-empty. Groups the finding, conventionally the
discipline name. \\
\code{severity} & string & Exactly one of \code{critical}, \code{warning},
\code{info}. \\
\code{message} & string & Non-empty. What the reader is told when a group is
missing. \\
\code{suggested\_fix} & string & Non-empty. What to do about it. \\
\code{requirement\_groups} & list of objects & Non-empty, with distinct
\code{label} values. Each is \code{\{"label": ..., "aliases": [...]\}}, the
aliases being a non-empty, duplicate-free list of the wordings that count as
mentioning the thing. \\
\end{longtable}

An alias may legitimately carry a trailing space --- \code{"iso~"} so that it
does not also match \code{isolation} --- so the emptiness check strips, but the
stored alias does not.

\file{profile\_catalog.py} validates all of this at import time, unknown keys
included. A typo is therefore an import-time error naming the file and the path
inside it (\code{nvh.json: rules[2].severity: must be one of critical, warning,
info}), never a rule that quietly stops firing.

Profile selection is \code{auto} by default: the catalog discipline of the
report decides, or failing that the identity patterns in the profile files,
tried in \code{detect\_priority} order with the first match kept.

\subsection{Semantic review}

An optional pass, requiring \env{REPORT\_LLM\_*} and a reachable Ollama. It asks
for structured JSON output and looks for three things:

\begin{itemize}
\item \code{semantic.scope\_result\_alignment} --- do the results answer the
stated scope?
\item \code{semantic.unsupported\_conclusion} --- is each conclusion grounded in
the report itself?
\item \code{semantic.internal\_contradiction} --- does the text contradict
itself?
\end{itemize}

Only findings whose quoted evidence can be verified \textbf{verbatim on the page
they claim} survive; anything the model invented or left unsupported is dropped
before it reaches a reviewer. With the LLM disabled the deterministic findings
stand alone.

\subsection{Human decisions and precision}

A finding is a hypothesis until a person rules on it.

\begin{shellblock}
POST /report-review/decisions
  {document_id, finding_key, rule_id,
   decision: open | confirmed | dismissed, note, reviewer}
\end{shellblock}

Decisions are stored in \code{report\_review\_decisions}, unique per
\code{(document\_id, finding\_key)}. \route{GET /report-review/rule-precision}
then reports the confirm rate per rule --- how often reviewers agreed with it.

A rule needs at least \code{MINIMUM\_PRECISION\_DECISIONS} (10) decided findings
before a rate is reported: below that, one reviewer dismissing two findings
would read as 0\,\% precision, which is noise rather than a measurement.

\subsection{Evidence and export}

\begin{description}
\item[\route{GET /documents/\{id\}/review-preview}] The report PDF with the
findings for one \code{rule\_id} highlighted on the requested page. Severity
picks the colour (critical red, warning amber, info blue).

\item[\route{GET /report-review/export}] A PDF review record for the requested
documents: the summary table, then each finding with its evidence, severity,
suggested fix and the human decision recorded against it.
\end{description}

\subsection{Revision comparison}

When two reports are selected, the review engine compares their findings rather
than their full text, and reports them as \textbf{new}, \textbf{resolved} and
\textbf{still open}. It answers ``what changed about the quality of this report
between revisions'', which is a different and more tractable question than a
full redline.

\section{Report comparison}

\code{ReportComparisonService} compares report \emph{content}, either pairwise
(\route{POST /report-comparison}) or across a set
(\route{POST /report-comparison/multi}, \code{reference} or \code{all\_pairs}
mode).

\begin{itemize}
\item Either side may be a document already in the database
(\code{document\_id}) or a one-off upload
(\route{POST /report-comparison/upload} returns an \code{upload\_token} valid
for a short window; the file is never ingested).
\item Chunks are aligned semantically and lexically; the result is split into
\textbf{similarities} and \textbf{differences}, each with a topic, a summary,
evidence from both sides and a confidence.
\item With \code{use\_llm} and a working provider, the LLM refines topics and
summaries. Without it the deterministic alignment stands.
\item \route{GET /report-comparison/\{id\}/pdf/\{side\}} returns the PDF with the
matched passages highlighted and numbered, and
\route{GET /report-comparison/\{id\}/viewer} opens a two-pane full-screen
viewer.
\end{itemize}

Results are cached per comparison id, which is what makes the viewer and the two
PDF endpoints cheap to open after the comparison itself has run.

\section{Report quality questions}

\code{ReportQualityService} answers a narrow class of question directly from the
document structure rather than from retrieval --- ``are the tables numbered
correctly?'', ``which figures are missing?''. It extracts caption occurrences
(\code{Tablo 3.2}, \code{Şekil 4}, \code{Resim 1}), checks the sequence, and
answers with the page each conclusion came from.
\code{DocumentIntelligenceService} routes matching chat questions here
automatically.

\section{Report drafting}

\code{ReportWriterService} (\route{POST /draft-report}) builds a report draft
from a title, an objective, free-form notes and an optional set of source
documents.

\begin{itemize}
\item Keywords are refined from the title, objective and notes; a retrieval query
is built from the result and run against the corpus.
\item The draft is composed from a cover block plus a body grounded in the
retrieved passages. \code{detail\_level} selects \code{quick} or
\code{detailed}.
\item With \env{REPORT\_LLM\_*} available the body is generated and then
\textbf{sanitized} against the cover block and the allowed terms; otherwise a
deterministic composition is used. Either way the response lists the sources it
drew on.
\item \route{POST /draft-report/pdf} renders the same payload as a PDF with a
cover table, a summary table and a Turkish-capable font.
\end{itemize}
